% Generated draft: requires mathematical and source review.
\documentclass[11pt]{article}
\usepackage[margin=25mm]{geometry}
\usepackage{fontspec}
\setmainfont{Latin Modern Roman}
\usepackage{amsmath,amssymb,mathtools,unicode-math}
\setmathfont{Latin Modern Math}
\usepackage{enumitem,needspace,xurl,hyperref}
\setlength{\parindent}{0pt}
\setlength{\parskip}{5pt}
\setlength{\emergencystretch}{3em}
\newcommand{\N}{\mathbb{N}}
\newcommand{\Z}{\mathbb{Z}}
\newcommand{\Q}{\mathbb{Q}}
\newcommand{\R}{\mathbb{R}}
\newcommand{\C}{\mathbb{C}}
\newcommand{\F}{\mathbb{F}}
\newcommand{\A}{\mathbb{A}}
\newcommand{\PP}{\mathbb{P}}
\newcommand{\E}{\mathbb{E}}
\newcommand{\eps}{\varepsilon}
\newcommand{\dd}{\,\mathrm{d}}
\begin{document}
\section*{UnsolvedMath ID: 30000043}
\textbf{Source problem code:} catalogue entry
\subsection{Definitions and mathematical statement}
\textbf{Question.} Does the canonical normalized tight frame associated with every wavelet frame always have wavelet structure? For a frame \((f_i)_{i\in I}\) in a Hilbert space \(H\), let \(Sf=\sum_i\langle f,f_i\rangle f_i\) be its positive invertible frame operator, with inner product linear in the first argument. Its canonical normalized tight frame is \((S^{-1/2}f_i)_i\), where the inverse square root is the positive one; normalized tight means Parseval:\[\sum_i|\langle f,S^{-1/2}f_i\rangle|^2=\|f\|^2\quad(f\in H).\] This is different from the canonical dual \((S^{-1}f_i)_i\). The question concerns retention of wavelet structure, not existence of the canonical Parseval frame.\par In the usual one-dimensional dyadic single-generator interpretation, put \(D f(x)=2^{1/2}f(2x)\), \(T^k f(x)=f(x-k)\), and suppose \((D^jT^k\psi)_{j,k\in\Z}\) is a frame for \(L^2(\R)\). The question then asks whether there is a generator \(\psi^{\mathrm{tight}}\in L^2(\R)\) such that\[S^{-1/2}D^jT^k\psi=D^jT^k\psi^{\mathrm{tight}}\quad\text{for every }j,k\in\Z.\] With this fixed indexing, \(\psi^{\mathrm{tight}}=S^{-1/2}\psi\), so the condition is equivalently \(S^{-1/2}D^jT^k\psi=D^jT^kS^{-1/2}\psi\) for all indices. The source does not specify the dimension, dilation, number of generators, or whether wavelet structure permits reindexing; this standard interpretation is therefore explicit and conditional, rather than a silent restriction of the general wording.\par The supplied literature triage checked 2026-08-21 labels the item open. It reports that canonical dual wavelet frames need not preserve wavelet structure in general, but that no theorem or counterexample settling the distinct canonical-tight question was found in the checked literature. A negative result about \(S^{-1}\) is not asserted here to settle the stated \(S^{-1/2}\) question. The assessment is dated and requires independent verification.
\subsection{Short English statement}
Applying the inverse square root of a wavelet frame operator always gives a Parseval frame. The question asks whether this particular frame is again generated by a wavelet with the same dilation and translation structure. It must be distinguished from the canonical dual-frame question.
\subsection{Sources}
\begin{itemize}\item Local UnsolvedMath registry, permanent ID 30000043, code catalogue entry; \emph{Wavelet Structure of Canonical Tight Frames}, Oberwolfach Reports Open Problems.\item \emph{Mini-Workshop: Wavelets and Frames (2004)}, Oberwolfach Reports, DOI 10.4171/owr/2004/10, pp. 479--544. \url{https://doi.org/10.4171/owr/2004/10}. Proposed year 2004; no proposer named.\item Supplied literature label \emph{OWR 10/2004}: \url{https://doi.org/10.4171/OWR/2004/10}.\item Supplied literature label \emph{Canonical dual wavelet frames}: \url{https://doi.org/10.1006/acha.2002.0381}.\item Supplied literature label \emph{Wavelet frame survey}: \url{https://arxiv.org/abs/1901.08024}.\item Literature triage checked 2026-08-21, status \texttt{open}; record updated 2026-09-16. Verify the dated assessment and cited claims independently.\end{itemize}
\end{document}
